%! Algebra 2 — Unit 1, Lesson 2 — Group Activity (Key)
\documentclass[10pt]{article}
\usepackage{algebra2-article}
\usepackage{algebra2-boxes}
\usepackage{algebra2-key}

\newcommand{\TallMath}[1]{$\displaystyle #1\rule[-1.4em]{0pt}{3.2em}$}

%=============================================================================
\begin{document}

\pageheader{Unit 1, Lesson 2}{Group Activity --- Key}
\namepartnerperiod

\vspace{0.15in}

% ─────────────────────────────────────────────
%  TIER R — Remediate
% ─────────────────────────────────────────────
\begin{tcolorbox}[colback=white,colframe=black!40,boxrule=0.7pt,arc=2mm,
    title={\bfseries Tier R — Remediate},fonttitle=\small,breakable]

\begin{enumerate}

    \item Solve $2x + 7 = 15$.

    \ansline{$2x = 15 - 7 = 8$}
    \ansline{$x = 4$}
    Answer: \ans{$x = 4$}

    \vspace{0.1in}

    \item Solve $3x - 1 > 8$; graph and identify notation.

    \ansline{$3x > 9 \Rightarrow x > 3$}
    Answer: \ans{B. \quad $x > 3$; \quad $(3, \infty)$}

    \begin{teachernote}
        Watch for students who use a closed dot at $3$ instead of open (strict inequality). Emphasize: $>$ means open dot, no endpoint included.
    \end{teachernote}

    \vspace{0.1in}

    \item Solve for $b$: $A = \dfrac{1}{2}bh$

    \ansline{$2A = bh$}
    Answer: \ans{$b = \dfrac{2A}{h}$}

    \vspace{0.1in}

    \item How many solutions does $4x + 2 = 4x - 3$ have?

    \ansline{$4x + 2 = 4x - 3 \Rightarrow 2 = -3$; contradiction}
    Answer: \ans{C. No solution}

    \begin{teachernote}
        Students may subtract $4x$ and get $2 = -3$ without recognizing this as a contradiction. Ask them: ``Is $2 = -3$ ever true?''
    \end{teachernote}

\end{enumerate}
\end{tcolorbox}

\vspace{0.2in}

% ─────────────────────────────────────────────
%  TIER A — Approaching Proficiency
% ─────────────────────────────────────────────
\begin{tcolorbox}[colback=white,colframe=black!40,boxrule=0.7pt,arc=2mm,
    title={\bfseries Tier A — Approaching Proficiency},fonttitle=\small,breakable]

\begin{enumerate}

    \item Solve $4(x - 3) = 2x + 6$.

    \ansline{$4x - 12 = 2x + 6$}
    \ansline{$2x = 18 \Rightarrow x = 9$}
    Answer: \ans{$x = 9$}

    \vspace{0.1in}

    \item Solve $-3x + 7 \leq 1$; write in interval notation.

    \ansline{$-3x \leq -6$}
    \ansline{Divide by $-3$; flip symbol: $x \geq 2$}
    Answer: \ans{$[2, \infty)$}

    \begin{teachernote}
        Most common error: forgetting to flip the inequality when dividing by $-3$.
    \end{teachernote}

    \vspace{0.1in}

    \item Solve system: $x + y = 5$, $\;2x - y = 4$.

    \ansline{From first equation: $x = 5 - y$}
    \ansline{Sub into second: $2(5-y) - y = 4 \Rightarrow 10 - 2y - y = 4$}
    \ansline{$-3y = -6 \Rightarrow y = 2$; then $x = 5 - 2 = 3$}
    Answer: \ans{$(3, 2)$}

    \vspace{0.1in}

    \item Solve $x^2 + 2x - 15 = 0$ by factoring.

    \ansline{Need two numbers: add to $2$, multiply to $-15$ $\Rightarrow$ $5$ and $-3$}
    \ansline{$(x + 5)(x - 3) = 0$}
    Answer: \ans{$x = -5$ or $x = 3$}

\end{enumerate}
\end{tcolorbox}

% ─────────────────────────────────────────────
%  TIER E — Extension
% ─────────────────────────────────────────────
\begin{tcolorbox}[colback=white,colframe=black!40,boxrule=0.7pt,arc=2mm,
    title={\bfseries Tier E — Extension},fonttitle=\small,breakable]

\begin{enumerate}

    \item Solve for $r$: $A = P(1 + rt)$

    \ansline{$A = P + Prt$}
    \ansline{$A - P = Prt$}
    Answer: \ans{$r = \dfrac{A - P}{Pt}$}

    \vspace{0.1in}

    \item Solve $\dfrac{1}{2}(x + 4) > 3x - 5$; interval notation.

    \ansline{$\tfrac{1}{2}x + 2 > 3x - 5$}
    \ansline{$7 > \tfrac{5}{2}x \Rightarrow x < \tfrac{14}{5}$}
    Answer: \ans{$\left(-\infty,\, \dfrac{14}{5}\right)$}

    \vspace{0.1in}

    \item Solve by elimination: $3x + 2y = 11$, $\;2x - 5y = -1$

    \ansline{Multiply eq.\ 1 by $5$, eq.\ 2 by $2$:}
    \ansline{$15x + 10y = 55$ \quad and \quad $4x - 10y = -2$}
    \ansline{Add: $19x = 53 \Rightarrow x = \dfrac{53}{19}$}
    \ansline{Sub back: $3\!\cdot\!\tfrac{53}{19} + 2y = 11 \Rightarrow 2y = 11 - \tfrac{159}{19} = \tfrac{50}{19} \Rightarrow y = \dfrac{25}{19}$}
    Answer: \ans{$\left(\dfrac{53}{19},\; \dfrac{25}{19}\right)$}

    \begin{teachernote}
        The non-integer solution is intentional --- reinforces that elimination works regardless of whether
        the answer is a nice number. Encourage Desmos verification: graph both lines and confirm the
        intersection. Students who get stuck on the fractions benefit from seeing that the method
        is correct even when arithmetic is messy.
    \end{teachernote}

    \vspace{0.1in}

    \item Solve $2x^2 - 5x - 3 = 0$ using the quadratic formula.

    \ansline{$\Delta = (-5)^2 - 4(2)(-3) = 25 + 24 = 49 > 0$; two real solutions}
    \ansline{$x = \dfrac{5 \pm \sqrt{49}}{4} = \dfrac{5 \pm 7}{4}$}
    \ansline{$x = \dfrac{12}{4} = 3$ \quad or \quad $x = \dfrac{-2}{4} = -\dfrac{1}{2}$}
    Answer: \ans{$x = 3$ or $x = -\dfrac{1}{2}$}

\end{enumerate}
\end{tcolorbox}

\end{document}
